% DSA5104 content adapted into the shared NUS Atlas cheatsheet layout.
% Two A4 portrait pages, four left-to-right columns per page, no global header/footer.
% Source text: user-provided DSA5104_A4_Helpsheet_v0.9_portrait_4col.tex.
% The source is treated as content; its document-level layout is replaced here.
\documentclass[a4paper]{article}
\usepackage[margin=4.5mm,left=5mm,right=5mm,noheadfoot]{geometry}
\usepackage{iftex}
\usepackage{amsmath,amssymb}
\ifPDFTeX
  \usepackage[T1]{fontenc}
  \IfFileExists{stix2.sty}{\usepackage{stix2}}{}
  \IfFileExists{sourcesanspro.sty}{\usepackage[default]{sourcesanspro}}{\usepackage{helvet}\renewcommand{\familydefault}{\sfdefault}}
  \IfFileExists{inconsolata.sty}{\usepackage{inconsolata}}{}
\else
  \usepackage{fontspec}
  \usepackage{unicode-math}
  \IfFontExistsTF{Source Sans 3}{\setmainfont{Source Sans 3}}{%
    \IfFontExistsTF{SourceSans3-Regular.otf}{\setmainfont{SourceSans3}[Extension=.otf,UprightFont=*-Regular,BoldFont=*-Bold,ItalicFont=*-It,BoldItalicFont=*-BoldIt]}{%
      \IfFileExists{sourcesanspro.sty}{\usepackage[default]{sourcesanspro}}{\setmainfont{lmsans10-regular.otf}[BoldFont=lmsans10-bold.otf]}}}
  \IfFontExistsTF{STIXTwoMath-Regular.otf}{\setmathfont{STIXTwoMath-Regular.otf}}{%
    \IfFontExistsTF{STIXTwoMath-Regular.otf}{\setmathfont{STIXTwoMath-Regular.otf}}{\setmathfont{latinmodern-math.otf}}}
  \IfFontExistsTF{JetBrains Mono}{\setmonofont{JetBrains Mono}[Scale=MatchLowercase]}{\setmonofont{lmmono10-regular.otf}[Scale=MatchLowercase]}
\fi
\usepackage{xcolor,array,enumitem,listings,multicol}
\definecolor{structure}{HTML}{10243A}
\definecolor{condition}{HTML}{704500}
\definecolor{conditionbg}{HTML}{FFF1CC}
\definecolor{warning}{HTML}{8A1C1C}
\definecolor{warningbg}{HTML}{FCE8E6}
\definecolor{example}{HTML}{174A6E}
\definecolor{examplebg}{HTML}{E8F1F8}
\definecolor{codebg}{HTML}{EFF3F7}
\lstdefinelanguage{SQLx}{
  morekeywords={SELECT,FROM,WHERE,GROUP,BY,HAVING,ORDER,ASC,DESC,DISTINCT,AS,JOIN,INNER,LEFT,RIGHT,FULL,OUTER,ON,USING,NATURAL,UNION,ALL,EXCEPT,INTERSECT,IN,NOT,EXISTS,WITH,RECURSIVE,INSERT,INTO,VALUES,UPDATE,SET,DELETE,CREATE,TABLE,VIEW,PROCEDURE,TRIGGER,BEFORE,AFTER,FOR,EACH,ROW,BEGIN,END,DECLARE,CALL,OUT,IF,THEN,ELSE,ELSEIF,SIGNAL,SQLSTATE,MESSAGE_TEXT,CASE,WHEN,NULL,IS,PRIMARY,KEY,FOREIGN,REFERENCES,CHECK,OVER,PARTITION,ROWS,RANGE,BETWEEN,PRECEDING,FOLLOWING,CURRENT,ROW,RANK,DENSE_RANK,ROW_NUMBER,SUM,COUNT,AVG,MIN,MAX,CUBE,ROLLUP,GROUPING,LIMIT},
  sensitive=false,
  morecomment=[l]{--},
  morestring=[b]'
}
\lstdefinelanguage{PHPx}{
  morekeywords={foreach,as,if,else,elseif,while,for,function,return,echo,true,false,null},
  sensitive=true,
  morecomment=[l]{//},
  morecomment=[l]{\#},
  morestring=[b]',
  morestring=[b]"
}
\lstset{
  language=SQLx,
  basicstyle=\ttfamily\fontsize{5.9}{6.7}\selectfont,
  keywordstyle=\color{structure},
  commentstyle=\color{example},
  stringstyle=\color{warning},
  backgroundcolor=\color{codebg},
  columns=fullflexible,
  keepspaces=true,
  showstringspaces=false,
  breaklines=true,
  breakatwhitespace=false,
  frame=none,
  aboveskip=1pt,
  belowskip=1pt,
  xleftmargin=1pt
}
\pagestyle{empty}
\setlength{\topskip}{0pt}
\setlength{\parindent}{0pt}
\setlength{\parskip}{1.3pt}
\setlength{\tabcolsep}{2pt}
\renewcommand{\arraystretch}{1.08}
\setlist[itemize]{leftmargin=8pt,label=\textbullet,nosep,topsep=1pt}
\setlength{\emergencystretch}{1em}
\newlength{\gutter}\setlength{\gutter}{2mm}
\setlength{\columnsep}{\gutter}\setlength{\multicolsep}{0pt}
\newcommand{\bodystyle}{%
  \fontsize{6.8}{8.6}\selectfont\raggedright
  \setlength{\abovedisplayskip}{2pt}\setlength{\belowdisplayskip}{2pt}%
  \setlength{\abovedisplayshortskip}{1pt}\setlength{\belowdisplayshortskip}{1pt}}
\newif\ifcolstart
\newcommand{\topic}[2]{\par\ifcolstart\colstartfalse\else\addvspace{5pt}\fi%
  \begingroup\setlength{\fboxsep}{0.8pt}%
  \noindent\colorbox{structure}{\parbox{\dimexpr\linewidth-2\fboxsep\relax}{%
    \color{white}\fontsize{8.1}{9.1}\selectfont\bfseries #1\enspace #2}}%
  \endgroup\par\nobreak\vspace{1pt}}
% Keep a topic together when the next column starts before its content ends.
\newenvironment{topicblock}
  {\par\begin{minipage}[t]{\linewidth}\bodystyle}
  {\end{minipage}\par}
\newcommand{\key}[1]{\par\textcolor{structure}{\textbf{#1}}\enspace}
\newcommand{\note}[4]{\par\begingroup\setlength{\fboxsep}{0.7pt}%
  \noindent\colorbox{#1}{\parbox{\dimexpr\linewidth-2\fboxsep\relax}{%
    \textcolor{#2}{\textbf{#3}}\enspace #4}}%
  \endgroup\par}
\newcommand{\cond}[1]{\note{conditionbg}{condition}{COND}{#1}}
\newcommand{\trap}[1]{\note{warningbg}{warning}{!}{#1}}
\newcommand{\ex}[1]{\note{examplebg}{example}{EX}{#1}}
\newcommand{\sql}[1]{\texttt{\detokenize{#1}}}
\newcommand{\code}[1]{\par\begingroup\setlength{\fboxsep}{1.2pt}%
  \noindent\colorbox{codebg}{\parbox{\dimexpr\linewidth-2\fboxsep\relax}{%
    \raggedright\color{structure}\ttfamily\fontsize{6.3}{7.7}\selectfont #1}}%
  \endgroup\par}

\begin{document}
\bodystyle
% PAGE 1: fill columns top-to-bottom, left-to-right; keep each topic intact.
\colstarttrue
\begin{multicols*}{4}
\begin{topicblock}
\topic{01}{QUERY STRUCTURE}

\key{Skeleton}
\begin{lstlisting}
SELECT [DISTINCT] expressions
FROM R [JOIN S ON condition]
[WHERE row_condition]
[GROUP BY keys]
[HAVING group_condition]
[ORDER BY expression [ASC|DESC]]
[LIMIT n];
\end{lstlisting}
\key{Logical processing order}
\code{FROM/JOIN -> WHERE -> GROUP BY}
\code{-> HAVING -> SELECT}
\code{-> ORDER BY -> LIMIT}
\trap{Window ORDER BY controls the window calculation, not the final output order. Add an outer ORDER BY when display order matters.}
\end{topicblock}

\begin{topicblock}
\topic{02}{NULL / AGGREGATION}

\sql{x = NULL} is wrong; use \sql{x IS NULL}. Comparisons with NULL yield UNKNOWN; \sql{WHERE} keeps only TRUE. Arithmetic with NULL usually returns NULL.
\key{Counts}
\sql{COUNT(*)} counts output rows; \sql{COUNT(x)} ignores NULL; \sql{COUNT(DISTINCT x)} counts unique non-NULL values.
\cond{Every selected nonaggregate expression must be grouped in standard SQL. \sql{WHERE} filters rows before grouping; \sql{HAVING} filters groups after aggregation.}
\trap{After a LEFT JOIN, COUNT(*) still counts a NULL-extended row; COUNT(B.id) ignores it.}
\end{topicblock}

\begin{topicblock}
\topic{03}{JOIN / SET / SUBQUERY}

\key{Which rows survive?}
\begin{itemize}
\item INNER: matched rows only.
\item LEFT: every left row; missing right values become NULL.
\item RIGHT: every right row.
\item FULL OUTER: unmatched rows from both sides; MySQL has no direct FULL OUTER JOIN syntax.
\end{itemize}
\key{ON vs WHERE}
ON decides which rows match; join type decides which unmatched rows survive; WHERE is the final row filter.
\begin{lstlisting}
SELECT A.* FROM A LEFT JOIN B
  ON A.id=B.id AND B.flag=1;
\end{lstlisting}
Putting a right-side condition in WHERE can remove NULL-extended rows and lose LEFT preservation. Put match rules in ON.
\key{Join traps}
NATURAL JOIN matches every same-named column, which can join unrelated fields; \sql{USING(k)} names shared keys explicitly. Anti-join: LEFT JOIN, then WHERE the right key IS NULL.
\key{Set operations}
UNION / INTERSECT / EXCEPT require union-compatible columns. They remove duplicates by default; UNION ALL preserves duplicates.
\key{Subqueries}
\begin{lstlisting}
WHERE x IN (SELECT ...)
WHERE EXISTS (SELECT 1 FROM S WHERE ...)
WHERE NOT EXISTS (...)
\end{lstlisting}
Correlated subqueries refer to outer-row aliases. For “for all”, test that there is no missing requirement (e.g. required set EXCEPT possessed set is empty).
\trap{NOT IN with any NULL in its subquery can yield UNKNOWN; NOT EXISTS is safer for anti-joins.}
\end{topicblock}

\begin{topicblock}
\topic{04}{DDL / DML}

\begin{lstlisting}
CREATE TABLE R (
  id INT PRIMARY KEY,
  x VARCHAR(20) NOT NULL,
  dept VARCHAR(20),
  CHECK (...),
  FOREIGN KEY (dept) REFERENCES D(dept)
);
INSERT INTO R(...) VALUES (...);
UPDATE R SET x=... WHERE ...;
DELETE FROM R WHERE ...;
\end{lstlisting}
\cond{Foreign-key actions include ON DELETE/UPDATE CASCADE, SET NULL, and RESTRICT/NO ACTION.}
\key{CTE}
\begin{lstlisting}
WITH X AS (SELECT ...)
SELECT ... FROM X;
\end{lstlisting}
\end{topicblock}

\begin{topicblock}
\topic{05}{RANKING / WINDOWS}

For scores 100, 100, 90:
\begin{center}
\begin{tabular}{@{}lccc@{}}
Function & 100 & 100 & 90\\\hline
\sql{ROW_NUMBER} & 1 & 2 & 3\\
RANK & 1 & 1 & 3\\
\sql{DENSE_RANK} & 1 & 1 & 2
\end{tabular}
\end{center}
\sql{ROW_NUMBER} is sequential and unique; ties need a deterministic tie-breaker. RANK leaves gaps after ties; \sql{DENSE_RANK} does not.
\begin{lstlisting}
RANK() OVER (
  PARTITION BY dept
  ORDER BY score DESC)
\end{lstlisting}
PARTITION BY restarts within each group; ORDER BY sets window order. Ranking on grouped results happens after grouping/aggregation.
\key{Window form}
\begin{lstlisting}
window_fn(expr) OVER (
  [PARTITION BY p] [ORDER BY o]
  [ROWS frame])
\end{lstlisting}
Windows calculate per row without collapsing rows. Aggregate windows include SUM/AVG/COUNT/MIN/MAX; ranking functions are nonaggregate.
\key{Frames}
\begin{lstlisting}
-- Running: first row through current
SUM(x) OVER (ORDER BY d ROWS BETWEEN
  UNBOUNDED PRECEDING AND CURRENT ROW)
-- Moving: previous 2 plus current
AVG(x) OVER (ORDER BY d ROWS BETWEEN
  2 PRECEDING AND CURRENT ROW)
-- Symmetric: previous, current, next
AVG(x) OVER (ORDER BY d ROWS BETWEEN
  1 PRECEDING AND 1 FOLLOWING)
\end{lstlisting}
\trap{ROWS counts physical rows. RANGE groups ORDER BY peers, so ties can enlarge the frame.}
\end{topicblock}

\begin{topicblock}
\topic{06}{RECURSIVE CTE}

Use for transitive closure, hierarchies, and paths.
\begin{lstlisting}
WITH RECURSIVE reach(a,b) AS (
  -- Anchor: base edges
  SELECT a,b FROM edge
  UNION
  -- Recursive term extends prior rows
  SELECT r.a,e.b FROM reach r
  JOIN edge e ON r.b=e.a
)
SELECT * FROM reach;
\end{lstlisting}
The anchor seeds rows; the recursive term refers to the CTE and repeatedly extends it.
\cond{UNION removes duplicates and can help closure/cycle control. UNION ALL avoids deduplication and needs a safe termination strategy.}
\end{topicblock}

\begin{topicblock}
\topic{07}{PROCEDURES / TRIGGERS}

\key{MySQL stored procedure}
\begin{lstlisting}
DELIMITER //
CREATE PROCEDURE p(IN p_id INT, OUT p_cnt INT)
BEGIN
  DECLARE v INT DEFAULT 0;
  SELECT COUNT(*) INTO v FROM T WHERE id=p_id;
  SET p_cnt=v;
END//
DELIMITER ;
CALL p(7,@n); SELECT @n;
\end{lstlisting}
IN receives a caller value; OUT returns one. Declare locals/conditions/handlers before executable statements. SELECT ... INTO stores a scalar result. DELIMITER is a client command so internal semicolons stay in the procedure body.
\key{Triggers}
Fire on INSERT, UPDATE, or DELETE. BEFORE validates or changes NEW; AFTER runs after success, e.g. for audit work.
\begin{lstlisting}
CREATE TRIGGER trg BEFORE UPDATE ON T
FOR EACH ROW
BEGIN
  IF NEW.score < 0 THEN
    SIGNAL SQLSTATE '45000'
      SET MESSAGE_TEXT='bad score';
  END IF;
END;
\end{lstlisting}
Row images: INSERT uses NEW; UPDATE uses OLD and NEW; DELETE uses OLD. SIGNAL raises a user error; a trigger error aborts the invoking statement.
\end{topicblock}

\begin{topicblock}
\topic{08}{PIVOT / CUBE / ROLLUP}

\key{Conditional aggregation}
Rows become grouping keys; CASE values become output columns.
\begin{lstlisting}
SELECT location,
 SUM(CASE WHEN type='THEFT' THEN 1 ELSE 0 END) theft,
 SUM(CASE WHEN type='BATTERY' THEN 1 ELSE 0 END) battery
FROM crime GROUP BY location;
\end{lstlisting}
For amounts, replace 1 with amount. COUNT(CASE WHEN condition THEN 1 END) also counts matches because COUNT ignores NULL.
\key{Grouping sets}
For dimensions (a,b,c):
\begin{itemize}
\item CUBE(a,b,c): every subset, $2^3=8$ sets.
\item ROLLUP(a,b,c): prefix hierarchy (abc, ab, a, empty), detail to subtotals to total.
\item GROUPING(a)=1 when a is omitted in a subtotal/total; use it to distinguish generated NULL from data NULL.
\end{itemize}
MySQL supports ROLLUP. CUBE/general grouping-set syntax is commonly shown with PostgreSQL in course material.
\end{topicblock}

\begin{topicblock}
\topic{09}{LAST-MINUTE TRAPS}

\begin{itemize}
\item Window functions preserve row count; GROUP BY collapses rows.
\item RANK has gaps; \sql{DENSE_RANK} has none; \sql{ROW_NUMBER} is unique.
\item ROWS BETWEEN n PRECEDING AND CURRENT ROW covers at most $n+1$ physical rows at an interior row.
\item Recursive CTE: anchor first; recursive term must reach a fixpoint or stop.
\item Trigger row images: INSERT=NEW, DELETE=OLD, UPDATE=both.
\item GROUPING() separates subtotal NULL from a real NULL value.
\end{itemize}
\end{topicblock}
\end{multicols*}
\newpage
% PAGE 2: continue in the same top-to-bottom, left-to-right order.
\colstarttrue
\begin{multicols*}{4}
\begin{topicblock}
\topic{10}{DBMS / RELATIONAL MODEL}

\key{Why a DBMS?}
Controls redundancy/inconsistency, difficult access, isolation, buried integrity rules, partial updates, concurrency, and authorization.
\key{Abstraction}
Physical: how data is stored. Logical: which data and relationships exist. View: a user-specific representation.
\cond{Physical data independence means changing physical storage without changing the logical schema or applications.}
\key{Schema / instance}
Schema is structure; instance is a snapshot of tuples. A relation is a set of tuples (order is irrelevant); a domain is a set of allowed atomic values. NULL denotes unknown/missing.
\key{Keys}
Superkey uniquely identifies tuples; candidate key is a minimal superkey; primary key is the selected candidate key. A foreign key references a key in another relation; NULL is allowed only if the column permits it.
\end{topicblock}

\begin{topicblock}
\topic{11}{RELATIONAL ALGEBRA}

\[
\begin{array}{@{}ll@{}}
\text{Rows} & \sigma_{condition}(R)\\
\text{Columns} & \pi_{columns}(R)\\
\text{All pairs} & R\times S\\
\text{Matched pairs} & R\bowtie_{condition}S\\
\text{Either set} & R\cup S\\
\text{R but not S} & R-S\\
\text{Rename} & \rho_X(R)\\
\text{Assign result} & X\leftarrow E
\end{array}
\]
\key{RA to SQL}
$\sigma$ = WHERE; $\pi$ = SELECT DISTINCT; $\bowtie$ = JOIN ... ON; $\cup$ = UNION; $-$ = EXCEPT/anti-join; $\rho$ = alias AS.
\trap{RA projection removes duplicates. SQL SELECT preserves them unless DISTINCT is requested.}
\end{topicblock}

\begin{topicblock}
\topic{12}{PHP / MYSQL SYNTAX}

\key{PHP patterns}
\begin{lstlisting}[language=PHPx]
$x                 // variable
$s = $a . $b;      // concatenate strings
foreach ($xs as $x) { ... }
$row["column"]     // associative-array field
$_GET["x"]; $_POST["x"];
$conn = mysqli_connect(...);
$r = mysqli_query($conn, $sql);
$rows = mysqli_fetch_all($r, MYSQLI_ASSOC);
if ($_SERVER["REQUEST_METHOD"] === "POST") { ... }
\end{lstlisting}
\key{HTML form / dropdown}
\begin{lstlisting}[language=PHPx]
<form method="post">
  <select name="x">
    <option value="...">...</option>
  </select>
</form>
\end{lstlisting}
\end{topicblock}

\begin{topicblock}
\topic{13}{HTTP / MYSQL RESULTS}

\key{Request flow}
Browser request $\to$ web server $\to$ PHP interpreter $\to$ MySQL. The result returns to PHP, which creates HTML/text for the browser.
\key{Query result}
\sql{mysqli_query()} returns a result handle/object, not a PHP row array. Fetch rows with \sql{mysqli_fetch_all($r, MYSQLI_ASSOC)}.
\key{Associative arrays}
\sql{key => value}; with \sql{MYSQLI_ASSOC}, database column names become keys, e.g. \sql{$row["ID"]}.
\key{GET / POST}
\sql{$_GET} comes from URL query parameters; \sql{$_POST} comes from the request body. \sql{$_SERVER["REQUEST_METHOD"]} identifies the method received by the page.
\cond{Use POST for form-body submissions; reading a variable from \sql{$_GET}/\sql{$_POST} is separate from validating its value.}
\end{topicblock}
\end{multicols*}
\end{document}
